\documentclass[11pt]{article}

\usepackage{tikz}
\usepackage{geometry}
\usepackage{graphicx}
\usepackage{color}
\usepackage{url}
\usepackage[export]{adjustbox}

\setlength{\parindent}{0em}
\setlength{\parskip}{1em}

\usepackage{geometry}
\geometry{margin=1in}

\renewcommand*\arraystretch{1.05}

\begin{document}

\footnotesize
\begin{tikzpicture}
  \node at (0,0) {\includegraphics[height=1.1in, keepaspectratio = true]{Block_I}};
  \node at (13.4, -0.05) {\vspace{1cm}\begin{tabular}{l}
      \\ \sc Daniel J Eck \\
       \it Associate Professor \\
       \it Department of Statistics \\
       University of Illinois \\
       605 E. Springfield Ave., \\
       Champaign, Illinois 61820 \\
       \texttt{dje13@illinois.edu} \\
       \end{tabular} };
\end{tikzpicture}

\normalsize

Dear Professors Doss, He, and Chopin,

We are pleased to submit our manuscript, ``The Aptitude-Achievement Paradox: A Theory of Rank Reversals under Selective Observation,'' for consideration in \textit{Journal of the Royal Statistical Society: Series~B (Statistical Methodology)}.

This paper establishes the probabilistic foundations of a general statistical framework for studying rank reversals under selective observation. Specifically, we formalize the aptitude-achievement paradox probability as a new inferential target for cross-context comparison, develop the first theoretical treatment of Full House Modeling (a novel modeling framework introduced by Yan et al.\ (2025, \textit{Annals of Applied Statistics}, vol.~19, no.~2, pp.~1778--1799)) including a characterization theorem establishing the uniqueness of its rank-preserving coupling under natural assumptions, and derive explicit methods for quantifying paradox probabilities.

We believe the manuscript is well suited to \textit{Series~B} because it develops new statistical methodology supported by rigorous probabilistic theory while maintaining a clear connection to practical statistical problems involving selective observation and cross-context comparison. Although motivated by examples from sports analytics and psychometrics, the proposed framework is broadly applicable to problems arising in education, labor economics, admissions, and other settings where differences in population structure, selection mechanisms, and measurement environments complicate comparisons across contexts.

Our manuscript makes the following primary contributions:

\begin{itemize}
    \item We formalize the aptitude-achievement paradox as an explicit inferential target for cross-context comparison and study its probability under selective observation.
    \item We develop the probabilistic foundations of Full House Modeling, including a characterization result showing that the induced latent-observed coupling is the unique rank-preserving monotone transformation consistent with the model assumptions.
    \item Under identical observed achievement distributions across contexts, we derive explicit and distribution-free paradox probabilities and characterize their dependence on population and sampling structure.
	\item Under heterogeneous achievement distributions, we derive general bounds and obtain explicit Gaussian representations with exponentially controlled remainder terms and consistent plug-in estimation procedures.
    \item We show theoretically and numerically that paradox probabilities depend primarily on ratios of eligible population sizes and sampling fractions rather than their absolute magnitudes.
\end{itemize}



We disclose that Claude and ChatGPT were used to assist with the writing of our manuscript.

This manuscript has not been published or submitted elsewhere, and all authors have approved its submission. We look forward to your feedback and appreciate your time and consideration.

Sincerely,

Shen Yan \\
Daniel J. Eck

\end{document}